% 第7回　コンセプト立案
\session{7}{2}{11月11日（水）}{コンセプト立案}{AIと発想法}{yes}

\goals{
  \item 発想法（異業種との掛け合わせ、前提の反転、極端化）をAIとの対話で回せる
  \item 量を出してから選ぶ、という発想の手順を実践できる
  \item コンセプトを「誰に、何を、どのように」の形で言語化できる
}

\fillin{狙うペルソナ（要約）}

\work[60mm]{ワーク1}{AIで30案を出す}
\begin{promptbox}[異業種と掛け合わせる]
自動車ディーラーの店舗に、図書館、ホテルのラウンジ、カフェ、保育園、スポーツジム、公園の要素を掛け合わせた店舗コンセプトを、それぞれ2案ずつ出してください。各案は「誰に／何を／どのように」の形式で、1案3行以内にしてください。ペルソナ：（確定したペルソナの要約）
\end{promptbox}
\begin{promptbox}[前提を反転する]
自動車ディーラーの店舗について「当たり前」とされていることを10個挙げ、それぞれを反転させた場合にどんな店舗になるかを1行ずつ示してください。
\end{promptbox}
\instr{極端化（「もし〜しかなかったら」「もし10倍だったら」）も自分たちで指示を考えて試す。}
\begin{wstabh}{colspec={Q[c,wd=22mm,m] Q[c,wd=12mm,m] X[l,m]},row{2-Z}={ht=9mm}}
  方法 & 案の数 & 気になった案（一言で） \\
  掛け合わせ & & \\ 反転 & & \\ 極端化 & & \\
\end{wstabh}
\fillin{極端化で使った自分たちの指示}

\work[70mm]{ワーク2}{5案に絞る（ペルソナの言葉で判定）}
\begin{wstabh}{colspec={Q[c,wd=6mm,m] X[2,l,m] X[3,l,m] Q[c,wd=14mm,m]},row{2-Z}={ht=13mm}}
  & 案 & ペルソナは「行きたい」と言うか（本人の言葉で） & 判定\newline ○△× \\
  1 & & & \\ 2 & & & \\ 3 & & & \\ 4 & & & \\ 5 & & & \\
\end{wstabh}

\work[70mm]{ワーク3}{コンセプト文を書く}
\begin{wstab}{colspec={Q[l,wd=30mm,m] X[l,m]},row{1-Z}={ht=13mm}}
  \hc{誰に} & \\
  \hc{何を} & \\
  \hc{どのように} & \\
  \hc{コンセプト文\newline（一文で）} & \\
  \hc{予備案} & \\
\end{wstab}

\work{発表}{各グループ1分。他グループの案で印象に残ったもの}
\writelines{3}

\begin{nexttime}
  \item コンセプトシート1枚（コンセプト文、狙うペルソナ、根拠、参考にした事例）
  \item コンセプトを実現する空間に必要な要素を、思いつく限り書き出す
\end{nexttime}
\twoboxes{根拠（企業の課題・ペルソナとの関係）}{参考にした事例（出典）}{32mm}
\answerbox[空間に必要な要素（思いつく限り）]{28mm}
\aimemo
